%% kapa: exam=1 round=30 year=2019 period=1 form=A subject=civil_law subject_ko=민법 seq=32 number=32
%% source: Q-Net 감정평가사 기출문제 (cst00309) artlSeq=5205205 file=extracted/1008/5205205/2019년 제30회 감정평가사 1차시험 문제지/2019년 제30회 감정평가사 1차시험 문제지/2019년도 제30회 감정평가사 1차 1교시 A형.hwp
%% conversion: kordoc 4.24.51 -> tools/gen1.py (auto); review=proofread-sonnet; flags=none
\begin{question}{32}
민법상 공동소유에 관한 설명으로 옳지 않은 것은? (다툼이 있으면 판례에 따름)
\begin{choices}
  \item 합의에 의한 공유물 분할의 경우, 공유자는 다른 공유자가 취득한 물건에 대하여 그 지분의 비율로 매도인과 동일한 담보책임이 있다.
  \item 공유자는 그 지분을 처분할 수 있고 공유물 전부를 지분의 비율로 사용, 수익할 수 있다.
  \item 공유자는 다른 공유자의 동의 없이 공유물을 처분하거나 변경할 수 있다.
  \item 공유물의 관리에 관한 사항은 공유자의 지분의 과반수로써 결정한다.
  \item 토지공유자 사이에서는 그 지분의 비율로 공유물의 관리비용 기타 의무를 부담한다.
\end{choices}
\end{question}
\answer{3}
